\rok{Јануар 2024.}{F}

Други практични тест из Алгоритама и структура података

Други практични тест одржан 10.01.2024.

\zadatak{1}{Завршавање имплементације BFS алгоритма}
Основна метода BFS алгоритма није потпуна. Допунити имплементацију ове методе која треба да омогући комплетно извршавање овог алгоритма.

\textbf{Додатне напомене за решавање задатка:}
\begin{itemize}
	\item Написати алгоритам.
	\item У Main методи креирати произвољан граф. Обезбедити начин за тестирање и проверу нове функционалности, односно позив BFS методе.
	\item У template-у је закоментарисана метода:
	\begin{Verbatim}[fontsize=\footnotesize, numbers=left, xleftmargin=10mm]
public void BFS(int startVertex)
	\end{Verbatim}
\end{itemize}



\zadatak{2}{Рачунање суме чворова који нису на позицији листова}
У оквиру стандардне структуре класе \texttt{BST} написати имплементацију методе која ће омогућити исписивање свих чворова који се налазе на позицији листова.

\textbf{Додатне напомене за решавање задатка:}
\begin{itemize}
	\item У оквиру класе \texttt{BST} креирати јавну методу са потписом:
	\begin{Verbatim}[fontsize=\footnotesize, numbers=left, xleftmargin=10mm]
public void sumNonLeaves()
	\end{Verbatim}
	\item Из претходно креиране јавне методе треба да се позива приватна метода са потписом:
	\begin{Verbatim}[fontsize=\footnotesize, numbers=left, xleftmargin=10mm]
private Node sumNonLeavesBST(Node curr)
	\end{Verbatim}
	\item Приватна метода треба да садржи имплементацију алгоритма којим ће се омогућити пролазак кроз стабло од корена ка листовима и њихово сумирање.
	\item Листом се класификује сваки чвор који нема ниједно дете.
	\item Може се креирати помоћна променљива на нивоу класе \texttt{BST} која ће помоћи у пребројавању.
	\item Алгоритам \textbf{ОБАВЕЗНО} писати у оквиру класе \texttt{BST}.
\end{itemize}

\textbf{Примери:}

\begin{center}
\begin{minipage}{\textwidth}
\centering
\begin{forest}
for tree={
	circle,
	draw,
	thick,
	fill=gray!20,
	minimum size=8mm,
	inner sep=1pt,
	s sep=8mm,
	l sep=12mm,
	edge={-, thick},
	font=\small
}
[8
	[4
		[2
			[, phantom]
			[3]
		]
		[5]
	]
	[9
		[, phantom]
		[11]
	]
]
\end{forest}

\vspace{0.3cm}
\textbf{Слика 1.} Пример BST-а.
\end{minipage}
\end{center}

\begin{itemize}
	\item \textbf{Улаз 1:} BST са слике 1
	\item \textbf{Излаз 1:} 23 (8 + 9 + 4 + 2)
\end{itemize}

\begin{center}
\begin{minipage}{\textwidth}
\centering
\begin{forest}
for tree={
	circle,
	draw,
	thick,
	fill=gray!20,
	minimum size=8mm,
	inner sep=1pt,
	s sep=8mm,
	l sep=12mm,
	edge={-, thick},
	font=\small
}
[10
	[4
		[2
			[3]
			[, phantom]
		]
		[6
			[5]
			[9]
		]
	]
	[20
		[12
			[, phantom]
			[14]
		]
		[29
			[27]
			[, phantom]
		]
	]
]
\end{forest}

\vspace{0.3cm}
\textbf{Слика 2.} Пример BST-а.
\end{minipage}
\end{center}

\begin{itemize}
	\item \textbf{Улаз 2:} BST са слике 2
	\item \textbf{Излаз 2:} 83 (10 + 4 + 20 + 2 + 6 + 12 + 29)
\end{itemize}



\zadatak{3}{Естимирана сличност секвенци након k карактера}
Написати алгоритам који рачуна естимирану сличност секвенци након прослеђених \texttt{k} карактера. Алгоритам као параметре добија секвенце и број \texttt{k} који означава почетак subsekvenci чију сличност треба естимирати. Алгоритам имплементирати уз претпоставку да је број \texttt{k} сигурно мањи од дужина секвенци које се естимирају.

\textbf{Додатне напомене за решавање задатка:}
\begin{itemize}
	\item Написати алгоритам.
	\item У Main методи обезбедити тестирање, позивом методе са параметрима који су дати у примерима.
	\item Користити \texttt{search} методу класе \texttt{KarpRabinAlgorithm} и класе \texttt{MinHash} и \texttt{Jaccard}, које су дате у шаблону.
	\item Имплементацију писати у оквиру закоментарисане методе:
	\begin{Verbatim}[fontsize=\footnotesize, numbers=left, xleftmargin=10mm]
public static double percentOfSimilarityAfterK(string str1, string str2, int k)
	\end{Verbatim}
	\item Приликом креирања \texttt{MinHash} објекта, за вредности параметара \texttt{universeSize} и \texttt{numHashFunctions} прослеђивати вредности 12 и 4.
	\item Параметар \texttt{q} функције \texttt{search} поставити на \texttt{substr.Length+1}.
	\item Дозвољено је изменити повратни тип методе \texttt{search} (нпр. уместо \texttt{void} имати \texttt{List<int>}).
\end{itemize}

\textbf{Напомене које могу бити од помоћи приликом решавања задатка:}
\begin{itemize}
	\item За конверзију string-а у низ карактера користити \texttt{ToCharArray()} по принципу: \texttt{str.ToCharArray()}.
	\item За конверзију низа карактера у string користити \texttt{new string(char[] arrayOfChars)} по принципу: \texttt{new string(charArray)}.
	\item Карактере није потребно експлицитно конвертовати у целобројни тип приликом убацивања у \texttt{List<uint>} структуру.
\end{itemize}

\textbf{Пример:}
\begin{itemize}
	\item \textbf{Улаз:} \texttt{"ACCCBBWMKCCCBCAC", "HRCCMKCPSAAKYCRBHVB", 10}
	\item \textbf{Излаз:} 0,2
	\item \textbf{Тумачење:} Естимирана сличност subsekvenci \texttt{CCBCAC} и \texttt{AKYCRBHVB} је 0,2.
\end{itemize}



\sablon{Шаблон другог теста јануар 2024. група F}{2023/Januar/GrupaF/AISP2023_Test2_GrupaF.linq}
